\origpage{47}{N. Robertson, P.D. Seymour / Journal of Combinatorial Theory, Series B 89 (2003) 43--76}
We have organized the paper backwards, to try to motivate the various steps better. Thus we first prove 1.3 assuming 3.1, then prove 3.1 assuming another statement 4.1, and so on.

\sectionhead{2. Structure relative to a tangle}

A \textit{separation} of a graph $G$ is a pair $(A, B)$ of subgraphs with $A \cup B = G$ and $E(A \cap B) = \emptyset$, and its \textit{order} is $|V(A \cap B)|$. A \textit{tangle of order $\theta$ in $G$}, where $\theta \ge 1$ is an integer, is a set $\mathscr{T}$ of separations of $G$, all of order $<\theta$, such that

\hspace*{12pt}1. For every separation $(A, B)$ of order $<\theta$, one of $(A, B)$, $(B, A)$ belongs to $\mathscr{T}$

\hspace*{12pt}2. If $(A_1, B_1)$, $(A_2, B_2)$, $(A_3, B_3) \in \mathscr{T}$ then $A_1 \cup A_2 \cup A_3 \ne G$, and

\hspace*{12pt}3. If $(A, B) \in \mathscr{T}$ then $V(A) \ne V(G)$.

We define $ord(\mathscr{T}) = \theta$.

A \textit{design} is a pair $(H, M)$ where $H$ \del{in} \fix{is}\fixnote{fix:design}{The word ``in'' stands for ``is''.} a graph and $M$ is a set of subsets of $V(H)$. A \textit{location} in $G$ is a set $\{(A_1, B_1), \ldots, (A_k, B_k)\}$ of separations of $G$ such that $A_i \subseteq B_j$ for all distinct $i, j$ with $1 \le i, j \le k$. If $\{(A_1, B_1), \ldots, (A_k, B_k)\}$ is a location then
\begin{equation*}
(G \cap B_1 \cap \cdots \cap B_k, \{V(A_i \cap B_i) \colon 1 \le i \le k\})
\end{equation*}
is a design, which we call the \textit{design of the location}.

Let $\theta \ge 1$ be an integer, and let $\mathscr{D}$ be a class of designs. We say that $\mathscr{D}$ is \textit{$\theta$-pervasive in a graph $G$} if for every subgraph $G'$ of $G$ and every tangle $\mathscr{T}$ in $G'$ of order $\ge \theta$ there is a location $\mathscr{L} \subseteq \mathscr{T}$ with its design in $\mathscr{D}$.

A \textit{tree-decomposition} of a graph $G$ is a pair $(T, \tau)$, where $T$ is a tree and for each $t \in V(T)$, $\tau(t)$ is a subgraph of $G$, such that

\hspace*{12pt}$\bullet$ $\cup(\tau(t) \colon t \in V(T)) = G$, and $E(\tau(t) \cap \tau(t')) = \emptyset$ for all distinct $t, t' \in V(T)$,

\hspace*{12pt}$\bullet$ if $t, t', t'' \in V(T)$ and $t'$ lies on the path of $T$ between $t$ and $t''$ then $\tau(t) \cap \tau(t'') \subseteq \tau(t')$.

If $(T, \tau)$ is a tree-decomposition of $G$ and $t_0 \in V(T)$, and $t_0$ has neighbours $t_1, \ldots, t_k \in V(T)$, then
\begin{equation*}
(\tau(t_0), \{V(\tau(t_0) \cap \tau(t_i)) \colon 1 \le i \le k\})
\end{equation*}
is a design, called the \textit{design of $t_0$ in $(T, \tau)$}.

Let $(H, M)$, $(H', M')$ be designs. We say that $(H', M')$ is an \textit{$n$-enlargement of $(H, M)$} if there exists $Z \subseteq V(H')$ such that

\hspace*{12pt}$\bullet$ $H$ is a subgraph of $H'$ and $V(H') \setminus V(H) \subseteq Z$,

\hspace*{12pt}$\bullet$ every edge of $H'$ is an edge of $H$,

\hspace*{12pt}$\bullet$ for every $X \in M'$ with $X \ne Z$, $X \cap V(H) \in M$,

\hspace*{12pt}$\bullet$ $n$ is an integer and $|Z| \le n$.

If $\mathscr{D}$ is a class of designs and $n \ge 0$ is an integer, the class of all $n$-enlargements of members of $\mathscr{D}$ is denoted by $\mathscr{D}^n$. The class of all designs $(H, M)$ with $|V(H)| \le n$ is denoted by $\mathscr{R}_n$. The following is implied by theorem 11.1 of [6].
